\documentclass[a4paper,11pt]{article}

% Francisation
\usepackage[utf8]{inputenc}
\usepackage[T1]{fontenc}
\usepackage{lmodern,textcomp}
\usepackage[frenchb]{babel}

\usepackage{amsfonts,amsmath,amssymb,mathrsfs}
\usepackage{geometry}
\usepackage{aeguill}
\usepackage{hyperref}
\usepackage{array}
\usepackage{multirow}
\usepackage{asymptote}
\usepackage{graphicx}
\usepackage{tabularx}
\usepackage{enumerate,tablists}
\usepackage{eurosym}
\usepackage{pst-solides3d}
\usepackage{multirow}
\usepackage{extsizes}
\usepackage{calc}
\usepackage{boites}
\usepackage{color,colortbl}
\usepackage{fancyhdr}
\usepackage{stmaryrd}
\usepackage{pifont}

\makeatletter
\let\Test@pr@shipout\pr@shipout
\let\Test@shipout\shipout
\makeatother
\usepackage{tikz,tkz-tab}
\makeatletter
\AtBeginDocument{%
  \let\pr@shipout\Test@pr@shipout
  \let\shipout\Test@shipout
}
\makeatother

\newtheorem{definition}{Définition}
\newtheorem{theoreme}{Théorème}

\newcommand{\C} {\ensuremath{\mathbb{C}}}
\newcommand{\R} {\ensuremath{\mathbb{R}}}
\newcommand{\Z} {\ensuremath{\mathbb{Z}}}
\newcommand{\N} {\ensuremath{\mathbb{N}}}
\newcommand{\ds} {\displaystyle}
\newcommand{\p}{\par}
\newcommand{\abs}[1]{\left |#1\right |}
\newcommand{\eq}[1]{\underset{#1}{\sim}}
\newcommand{\tend}[1]{\underset{#1}{\longrightarrow}}
\newcommand{\norme}[1]{\left\Vert #1\right\Vert}
\newcommand{\V}[1]{\overrightarrow{#1}}

\definecolor{gris1}{gray}{0.85}
\definecolor{gris2}{gray}{0.65}

\setlength{\textwidth} {19cm}
\setlength{\textheight} {26cm}
\addtolength{\topmargin} {-2.5cm}
\setlength{\oddsidemargin} {-1.5cm}
\setlength{\evensidemargin} {-1.5cm}

\pagestyle{fancy}
\fancyfoot[R]{\today}\renewcommand\headrulewidth{1pt}
\fancyhead[L]{ESTP / Bachelor 2 / \the\year}
\fancyhead[R]{Composition de mathématiques}
\renewcommand\footrulewidth{1pt}

\begin{document}
\begin{center}\bf {\Large\underline{Composition de MATHÉMATIQUES  (3h)}}\end{center}
\hfill\break

\vspace{2em}

\underline {\ding{110}\, \bf{Exercice 1 : Système différentiel couplé } \bf{(10 points)}}
\medskip

On cherche les triplets de fonctions $\left(x_1, x_2, x_3\right) \in \mathcal{C}^1(\mathbb{R})$ vérifiant, pour tout $t \in \mathbb{R}$,
\[
\left\{\begin{array}{l}
x_1^{\prime}(t)=x_1(t)+x_3(t), \\
x_2^{\prime}(t)=-x_1(t)+2 x_2(t)+x_3(t), \\
x_3^{\prime}(t)=2 x_1(t)-2 x_2(t).
\end{array}\right.
\]
On pose, pour $t \in \mathbb{R}$,

$$
X(t)=\left(\begin{array}{l}
x_1(t) \\
x_2(t) \\
x_3(t)
\end{array}\right) .
$$


\begin{enumerate}
    \item Écrire la matrice $A$ du système différentiel telle que $X^{\prime}(t)=A X(t)$.
    \item Diagonaliser la matrice $A$.
    \item Soit $\left(u_1, u_2, u_3\right)$ une base de vecteurs propres pour $A$ et $P$ la matrice de passage de la base canonique de $\mathbb{R}^3$ à la base $(u_1, u_2, u_3)$. On pose, pour tout $t \in \mathbb{R}$, $Y(t)=P^{-1} X(t)$. Quelle équation relie $Y^{\prime}$ et $Y$ ?
    \item Résoudre cette équation, puis en déduire les solutions du système différentiel initial.
\end{enumerate}

\hfill\break
\hrule
\hfill\break
\vspace{1em}

\underline {\ding{110}\, \bf{Exercice 2 : Chaîne de Markov et diagonalisation } \bf{(10 points)}}
\medskip

On modélise la circulation d'étudiants entre trois zones d'un campus :
\[
\text{B}=\text{Bibliothèque},\quad
\text{C}=\text{Cafétéria},\quad
\text{S}=\text{Salle informatique}.
\]
À chaque pas de temps (30 minutes), un étudiant change de zone selon les probabilités suivantes :
\begin{itemize}
    \item Depuis B : 50\% restent en B, 30\% vont en C, 20\% vont en S.
    \item Depuis C : 20\% vont en B, 60\% restent en C, 20\% vont en S.
    \item Depuis S : 30\% vont en B, 20\% vont en C, 50\% restent en S.
\end{itemize}

On note $P$ la matrice de transition (lignes dans l'ordre B, C, S), et $\pi_n$ le vecteur ligne des proportions à l'instant $n$.

\begin{enumerate}
    \item Écrire la matrice $P$. Vérifier que $P$ est stochastique. (1 point)
    \item Déterminer le polynôme caractéristique de $P$, puis ses valeurs propres. (2 points)
    \item Montrer que $P$ est diagonalisable et déterminer une matrice $Q$ inversible et une matrice diagonale $D$ telles que $P=QDQ^{-1}$. (3 points)
    \item En déduire une expression de $P^n$ pour tout $n\geq 1$. (2 points)
    \item À l'instant initial, la répartition est $\pi_0=(1,0,0)$. Exprimer $\pi_n$ puis déterminer $\displaystyle\lim_{n\to +\infty}\pi_n$. Interpréter ce résultat. (2 points)
\end{enumerate}

\hfill\break
\hrule
\hfill\break
\newpage

\underline {\ding{110}\, \bf{Exercice 3 : Matrice à paramètre et diagonalisation } \bf{(5 points)}}\\
\medskip
Soit $m$ un nombre réel et $f$ l'endomorphisme de $\mathbb{R}^3$ dont la matrice dans la base canonique est

$$
A=\left(\begin{array}{rcc}
1 & 0 & 1 \\
-1 & 2 & 1 \\
2-m & m-2 & m
\end{array}\right) .
$$

\begin{enumerate}
    \item Quelles sont les valeurs propres de $f$ ?
    \item Pour quelles valeurs de $m$ l'endomorphisme est-il diagonalisable ?
    \item On suppose $m=2$. Calculer $A^k$ pour tout $k \in \mathbb{N}$.
\end{enumerate}

\end{document}
